\documentclass[12pt,a4paper]{report}

\usepackage[T1]{fontenc}
\usepackage[utf8]{inputenc}
\usepackage{lmodern}
\usepackage[top=1in,bottom=1in,left=1.25in,right=1in]{geometry}
\usepackage{setspace}
\usepackage{graphicx}
\graphicspath{{figures/}}
\usepackage{booktabs}
\usepackage{array}
\usepackage{tabularx}
\usepackage{longtable}
\usepackage{multirow}
\usepackage{amsmath}
\usepackage{amssymb}
\usepackage{xcolor}
\usepackage{listings}
\usepackage[hidelinks]{hyperref}
\usepackage{fancyhdr}
\usepackage{titlesec}
\usepackage{caption}
\usepackage{subcaption}
\usepackage{float}
\usepackage{tikz}
\usepackage{pgfplots}
\pgfplotsset{compat=1.18}
\usetikzlibrary{shapes.geometric,arrows.meta,positioning,fit,backgrounds,calc}
\usepackage[style=ieee,backend=biber]{biblatex}
\addbibresource{references.bib}

\pagestyle{fancy}
\fancyhf{}
\fancyhead[L]{\small LeBlanc: Automated Prompt Enrichment and Iterative Repair}
\fancyhead[R]{\small \thepage}
\renewcommand{\headrulewidth}{0.4pt}

\titleformat{\chapter}[block]
  {\large\bfseries\centering}{Chapter \thechapter:}{0.5em}{}
\titlespacing*{\chapter}{0pt}{20pt}{10pt}

\onehalfspacing

% =========================================================
\begin{document}
\setcounter{page}{8}   % front matter occupies ~7 pages

% =========================================================
% CHAPTER 1 -- INTRODUCTION  (pages 6-9 approx)
% =========================================================
\chapter{Introduction}
\label{ch:introduction}

\section{Background and Motivation}
\label{sec:background}

The use of large language models (LLMs) as coding assistants has become
standard practice in professional software development. Tools such as GitHub
Copilot, ChatGPT, and Google Gemini generate syntactically correct, functionally
plausible code within seconds of receiving a natural-language prompt. Developer
surveys indicate that a majority of professional programmers now use at least one
AI-assisted coding tool regularly, and this proportion is growing.

This adoption creates a security problem that is less visible but well-documented.
LLMs are trained primarily to maximize functional correctness---to produce code
that runs and satisfies the stated requirement. They learn from public repositories
on platforms such as GitHub, where a substantial fraction of existing code contains
known security weaknesses. The model learns these patterns along with everything
else, and it reproduces them when asked to write similar code.

The consequence in practice is straightforward. A developer who asks an LLM to
``write a Flask login endpoint with MySQL'' receives working code that is, with
measurable probability, also exploitable. SQL injection parameters go unsanitised.
Password comparison happens in plaintext. Session tokens are returned without
appropriate flags. The developer sees the code run correctly against a test
database, ships it, and the vulnerability ships with it.

The 2022 study by Pearce et al.~\cite{pearce2022asleep} quantified this problem
against GitHub Copilot: 40\% of the generated code samples they evaluated
contained at least one security vulnerability, with a vulnerability rate of 63.4\%
across the LLMSecEval prompt set. Subsequent studies on ChatGPT and other models
found comparable or worse rates~\cite{khoury2023secure,tony2023llmseceval}. The
problem is not specific to one model or one era of models---it appears to be a
structural consequence of training on general-purpose code repositories without
security-specific supervision.

\section{Problem Statement}
\label{sec:problem}

Three partial solutions exist in the literature, each addressing one aspect of
the problem but not the whole.

The first class of solution modifies the training process: fine-tuning on
secure code pairs~\cite{hexacoder2024}, preference optimisation toward secure
outputs~\cite{lpo2024}, or constrained decoding at inference
time~\cite{codeguard2024}. These approaches can be effective but require access
to model weights or specialised training infrastructure. They cannot be applied
to closed-source commercial models, and they require a new training run for each
new model version.

The second class of solution is evaluation-only: run a set of prompts through
one or more models, scan the outputs, and report a vulnerability rate. The
SecurityEval~\cite{securityeval}, LLMSecEval~\cite{tony2023llmseceval}, and
CodeSecEval~\cite{codeseceval2023} benchmarks all fall into this category. These
studies produce useful comparative data but stop at detection. They do not attempt
to reduce vulnerability rates; they measure them.

The third class addresses the enrichment side: inject CWE-specific warnings into
the prompt before sending it to the model, on the hypothesis that explicitly
naming the vulnerability class the model should avoid will reduce its frequency.
The SecuCoGen study~\cite{tihanyi2023secucogenicse} demonstrated this approach
on GPT-3.5 and found substantial reduction. However, the study used a single
model, a fixed dataset of 180 samples, and did not include any post-generation
repair mechanism.

The gap these three classes leave is: a single deployable system that (a) enriches
prompts with CWE warnings automatically, (b) scans generated code with real static
analysis tools, and (c) drives an iterative repair loop until the scanner
returns clean, running all of this simultaneously across multiple models for
direct comparison. No such system exists in the literature.

\section{Proposed System: LeBlanc}
\label{sec:proposed}

LeBlanc fills this gap. The name refers to the Le Blanc error, a classic
software quality concept in which a defect introduced early in development is
significantly cheaper to fix at that point than after deployment---a principle
that applies directly to security vulnerabilities in generated code.

The system is a three-engine pipeline:

\begin{itemize}
  \item \textbf{Engine A --- Prompt Enrichment.} A rule-based local system that
        parses any coding prompt for security-relevant keywords and injects
        structured CWE warnings before the LLM call. No LLM is involved in this
        step; the enrichment is deterministic and reproducible.

  \item \textbf{Engine B --- Static Analysis.} Once the LLM returns code,
        Engine~B runs Semgrep and Bandit in parallel via subprocess calls and
        maps their raw JSON output to structured findings: CWE identifier,
        severity, line number, and vulnerability category.

  \item \textbf{Engine C --- Iterative Repair.} When Engine~B finds
        vulnerabilities, Engine~C builds a repair prompt from the scanner
        findings and resubmits it to the originating LLM. The patched output is
        immediately re-scanned. The loop runs for up to three iterations or until
        the scanner returns clean.
\end{itemize}

The pipeline operates against multiple LLMs simultaneously. For any given prompt,
results for all models and all modes (plain, enriched, enriched with repair) are
collected, scanned, and stored in a single database run. The frontend displays
a side-by-side comparison of vulnerability rates, scan findings, and repair
convergence for each model.

\section{Research Questions}
\label{sec:rqs}

LeBlanc is designed to answer five research questions from a single automated
data collection run:

\begin{description}
  \item[RQ1 --- Enrichment Effect.] Does automated CWE-aware prompt enrichment
        reduce vulnerability rates, and by how much per model?
  \item[RQ2 --- Model Comparison.] Which model produces the safest code out of
        the box? Does the ranking change after enrichment?
  \item[RQ3 --- Repair Effectiveness.] How many repair iterations does each model
        need to reach clean code? Is there a ceiling effect where some
        vulnerability types never converge?
  \item[RQ4 --- CWE Difficulty.] Which CWE categories are hardest for LLMs to
        avoid and hardest to self-repair?
  \item[RQ5 --- Generational Gap.] How does vulnerability rate and repair
        convergence differ between small older models (Llama 3.1-8B) and
        contemporary frontier models (Command-A, Gemini 2.5 Flash)?
\end{description}

\section{Scope and Limitations}
\label{sec:scope}

The current evaluation covers 33 prompts from the LLMSecEval benchmark, three
models, and up to three repair iterations per run. The dataset and model set are
intentionally bounded to allow one complete reproducible run within semester
constraints. The LLMSecEval prompts are Python-only; Engine B is correspondingly
configured for Python static analysis rules.

The Llama 3.1-8B results should be read carefully. An 8-billion parameter model
represents an earlier, smaller generation of open-weight LLMs. Its behaviour under
prompt enrichment---where adding CWE warnings actually increased vulnerability
rate---is consistent with published findings that small models degrade under
complex instruction sets rather than improving. These results are included for
completeness and as a motivating datapoint for the RQ5 generational gap analysis.
They should not be taken as evidence that prompt enrichment is counterproductive.
Future work will repeat the evaluation on Llama 3.3-70B, the direct
next-generation successor in the same model family.

\section{Contributions}
\label{sec:contributions}

This project makes three concrete contributions:

\begin{enumerate}
  \item \textbf{A working pipeline system.} LeBlanc is a fully implemented,
        runnable web application that takes a coding prompt, runs it through all
        three engines across multiple LLMs in parallel, and displays structured
        results. The source code and database schema are reproducible.

  \item \textbf{An empirical dataset.} Every pipeline run logs the original
        prompt, enriched prompt, raw LLM output, extracted code, Semgrep and
        Bandit findings, and every repair iteration to a SQLite database. The
        collected dataset covers 321 total runs across three models, three modes,
        and 33 prompts.

  \item \textbf{Comparative results across five metrics.} Vulnerability rate,
        relative improvement over the Copilot baseline, prompt enrichment gain,
        repair convergence rate, and mean iterations to fix are all reported per
        model and per CWE category against a well-known external baseline.
\end{enumerate}

\section{Report Organisation}
\label{sec:organisation}

Chapter~\ref{ch:literature} reviews the eight prior works that directly motivate
LeBlanc's design and identifies the specific gap each one leaves.
Chapter~\ref{ch:system} describes the system architecture, the three engines, and
the implementation choices.
Chapter~\ref{ch:methodology} defines the evaluation methodology: dataset
selection, metrics definitions, and experimental setup.
Chapter~\ref{ch:results} presents the results for all five research questions.
Chapter~\ref{ch:discussion} interprets the findings, discusses limitations, and
addresses the Llama 3.1-8B degradation result specifically.
Chapter~\ref{ch:conclusion} summarises the work and outlines future directions,
including the VS~Code extension and the planned larger-model evaluation.

% =========================================================
% CHAPTER 2 -- LITERATURE REVIEW  (pages 10-15 approx)
% =========================================================
\chapter{Literature Review}
\label{ch:literature}

This chapter reviews the eight papers that directly informed LeBlanc's design.
Each is summarised on its contribution and the gap it leaves. The gaps are not
criticisms---each paper is well-executed within its stated scope. They define the
boundary of existing knowledge that LeBlanc is designed to extend.

\section{Overview of the Review}
\label{sec:litoverview}

The review covers three themes: (1) vulnerability measurement studies that
establish baseline rates, (2) enrichment and fine-tuning approaches that attempt
to reduce rates at generation time, and (3) repair and feedback approaches that
attempt to fix vulnerabilities after generation. LeBlanc sits at the intersection
of all three.

Table~\ref{tab:litgap} summarises each paper against the capabilities present in
LeBlanc's design.

\begin{table}[H]
\centering
\caption{Prior work and the gaps addressed by LeBlanc}
\label{tab:litgap}
\small
\begin{tabularx}{\textwidth}{@{}lXXX@{}}
\toprule
\textbf{Paper} & \textbf{Core Contribution} & \textbf{Gap Left} & \textbf{LeBlanc's Answer} \\
\midrule
Pearce et al.\ (2022)~\cite{pearce2022asleep}
  & Copilot security audit; 63.4\% VR baseline
  & Single model; no repair; observational only
  & Multi-model; automated repair; active pipeline \\[6pt]

Tony et al.\ (2023)~\cite{tony2023llmseceval}
  & LLMSecEval dataset; 120 NL prompts
  & Single-turn evaluation; no enrichment or repair
  & Enrichment + repair on LLMSecEval prompts \\[6pt]

Khoury et al.\ (2023)~\cite{khoury2023secure}
  & ChatGPT security evaluation across 21 scenarios
  & Single model; no automated enrichment or repair
  & Multi-model parallel evaluation \\[6pt]

Tihanyi et al.\ (2024)~\cite{tihanyi2023secucogenicse}
  & CWE-aware prompting; GPT-3.5 VR from 31\% to 86\%
  & Single model; 180 samples; no repair loop
  & Generalises enrichment to multiple models + adds repair \\[6pt]

Wang et al.\ (2024)~\cite{codeseceval2023}
  & 44-CWE benchmark; automated evaluation
  & No enrichment engine; no iterative repair
  & All three modes in one system \\[6pt]

He \& Vechev (2024)~\cite{codeguard2024}
  & Constrained decoding for secure generation
  & Requires model internals; not API-accessible
  & API-only; works on any black-box LLM \\[6pt]

Siddiq et al.\ (2024)~\cite{hexacoder2024}
  & Oracle-guided secure fine-tuning (HexaCoder)
  & Fine-tune only; unusable on closed models
  & No fine-tuning; prompt-side intervention \\[6pt]

Zhou et al.\ (2024)~\cite{lpo2024}
  & Preference optimisation for secure generation
  & Requires retraining; not deployable on closed models
  & Fully deployment-agnostic; API-level only \\
\bottomrule
\end{tabularx}
\end{table}

\section{Vulnerability Measurement Studies}
\label{sec:litbaseline}

\subsection{Pearce et al.\ (2022) --- Asleep at the Keyboard}
\label{ssec:pearce}

Pearce et al.~\cite{pearce2022asleep} conducted the first systematic security
audit of GitHub Copilot's code generation. They constructed 89 scenarios spanning
CWE categories including injection, memory safety, and cryptographic weaknesses,
and evaluated Copilot's outputs using manual expert review and automated scanning.
The study found that approximately 40\% of all generated code samples contained at
least one security vulnerability, and that certain CWE categories---particularly
CWE-89 (SQL Injection) and CWE-22 (Path Traversal)---appeared at high rates
across multiple scenarios.

The headline figure from the LLMSecEval subset is 63.4\%: across 82 prompts,
52 produced vulnerable output. This number has become the standard external
baseline for subsequent LLM security evaluation work, and LeBlanc uses it as the
primary comparison point throughout the results chapter.

The study's limitation from LeBlanc's perspective is that it is observational
and single-model. It measures vulnerability rates but proposes no mechanism to
reduce them. There is no enrichment, no repair loop, and no comparison between
models. It establishes the problem but does not attack it.

\subsection{Tony et al.\ (2023) --- LLMSecEval}
\label{ssec:tony}

Tony et al.~\cite{tony2023llmseceval} released the LLMSecEval dataset: 120
natural-language prompts designed to elicit code containing common security
weaknesses, drawn from the CWE Top 25. Each prompt is annotated with its target
CWE and a ``vulnerable by design'' flag. The dataset covers Python scenarios for
weaknesses including SQL injection, path traversal, hardcoded credentials, and
weak cryptography.

LeBlanc uses a 33-prompt subset of LLMSecEval as its evaluation dataset. The
prompts are functional application requests---``write a Flask login endpoint,''
``build a file upload handler''---rather than explicitly asking the model to write
insecure code, which could trigger safety filters in commercial models.

Tony et al.'s own evaluation is single-turn: they send each prompt to several
models and measure how often the output contains the target CWE. No enrichment,
repair, or multi-turn analysis is performed. LeBlanc extends this by running each
LLMSecEval prompt through all three pipeline modes.

\subsection{Khoury et al.\ (2023) --- How Secure Is Code Generated by ChatGPT?}
\label{ssec:khoury}

Khoury et al.~\cite{khoury2023secure} evaluated ChatGPT across 21 scenarios
covering buffer overflows, SQL injection, cryptographic weaknesses, and other
common CWE types. They found that approximately 16 of the 21 scenarios produced
at least one piece of code containing a security vulnerability on the first
request, although requesting revision sometimes caused ChatGPT to produce a
secure version.

The manual nature of their repair---a human explicitly requests a
revision---distinguishes their work from LeBlanc's automated repair loop.
Their observation that revision requests can be effective is empirical support
for the premise behind Engine~C: the same LLM that produces vulnerable code can
often fix it when given specific feedback about what is wrong.

\section{Enrichment and Fine-Tuning Approaches}
\label{sec:litenrich}

\subsection{Tihanyi et al.\ (2024) --- SecuCoGen}
\label{ssec:secucogen}

Tihanyi et al.~\cite{tihanyi2023secucogenicse} is the most directly relevant
prior work to LeBlanc's Engine~A. They constructed a 180-sample Python dataset
of security-sensitive prompts and evaluated GPT-3.5-Turbo under two conditions:
standard prompts and CWE-enriched prompts that explicitly named the vulnerability
the model should avoid. For standard prompts, GPT-3.5 produced vulnerable code
in 31\% of cases. For CWE-enriched prompts, this dropped to 14\%---a reduction
of over 50\%.

This result directly motivates LeBlanc's design. Engine~A automates and
generalises the SecuCoGen enrichment strategy: instead of manually constructing
CWE-enriched prompts, Engine~A applies a keyword-to-CWE mapping to any
arbitrary prompt and injects the warnings automatically.

The limitations SecuCoGen leaves are: single-model (GPT-3.5 only), fixed
Python-only dataset, no post-generation repair, and no multi-model comparison.
LeBlanc addresses all four.

\subsection{Siddiq et al.\ (2024) --- HexaCoder}
\label{ssec:hexacoder}

HexaCoder~\cite{hexacoder2024} takes a different approach to reducing
vulnerability rates: oracle-guided fine-tuning. The system generates a large
training set of (vulnerable, secure) code pairs using an ``oracle''---a
combination of scanner feedback and expert review---and fine-tunes a base LLM on
this dataset. The resulting model generates substantially safer code than the
baseline.

The limitation relevant to LeBlanc is deployment: fine-tuning requires access
to model weights. It cannot be applied to closed-source commercial APIs such as
Gemini or Command-A. It also requires a significant training pipeline and compute
resources, making it inaccessible to most practitioners who would benefit from
safer code generation.

LeBlanc's design explicitly rejects any fine-tuning dependency. All three engines
operate at the API level. The system works against any model that can be accessed
via a standard chat completion API, regardless of whether the weights are
available.

\subsection{Zhou et al.\ (2024) --- LPO}
\label{ssec:lpo}

Zhou et al.~\cite{lpo2024} propose Language Model Preference Optimisation (LPO),
a reinforcement-learning-from-human-feedback variant applied to the security
domain. They train a reward model that scores generated code on security criteria
and fine-tune the base LLM to maximise that reward. The approach produces models
with improved security properties compared to standard instruction-tuned baselines.

The same deployment limitation applies as for HexaCoder: LPO requires full
access to model weights and training infrastructure. It cannot be retrofitted to
an existing closed-source API deployment.

\section{Constrained Decoding and Hybrid Methods}
\label{sec:litconstrained}

\subsection{He and Vechev (2024) --- CODEGUARD+}
\label{ssec:codeguard}

He and Vechev~\cite{codeguard2024} present CODEGUARD+, a constrained decoding
approach that modifies the token sampling process during generation to steer the
model away from patterns associated with known vulnerabilities. The constraints
are derived from a learned security classifier that operates on partial token
sequences.

CODEGUARD+ is technically sophisticated and demonstrably effective. However, it
requires direct access to the model's logits at each decoding step, which is only
possible when running the model locally. It cannot be applied to API-served models
where only the final generated text is returned. This rules out its application
to Gemini, Command-A, or any commercial model.

\section{Evaluation Benchmarks}
\label{sec:litbenchmark}

\subsection{Wang et al.\ (2024) --- CodeSecEval}
\label{ssec:codeseceval}

Wang et al.~\cite{codeseceval2023} developed CodeSecEval, a benchmark covering
44 CWE types with 180 samples. The benchmark provides both a generation task
(produce code satisfying a description) and a repair task (fix a provided
vulnerable snippet). They evaluate several models on both tasks and report
per-CWE difficulty rankings.

CodeSecEval is the broadest existing benchmark by CWE coverage, and its per-CWE
difficulty rankings directly inform LeBlanc's RQ4 (CWE difficulty). The key gap
is that CodeSecEval's repair task is static---a fixed vulnerable snippet is
provided to the model as input, and the model produces one repair attempt. There
is no iterative loop, no re-scanning after repair, and no automated detection of
whether the repair succeeded. LeBlanc's Engine~C closes this gap.

\section{Summary of Gaps}
\label{sec:litgapsummary}

Reviewing the eight papers, five consistent gaps emerge:

\begin{enumerate}
  \item \textbf{Single-model evaluation.} Six of the eight papers evaluate only
        one model. Cross-model comparative data under identical conditions does
        not exist prior to this work.

  \item \textbf{Single-turn evaluation.} Seven of the eight papers test generation
        once and stop. No study measures iterative repair under automated
        re-scanning.

  \item \textbf{No closed-loop enrichment.} SecuCoGen demonstrates that
        CWE-aware prompting works but requires manual prompt construction. No
        prior system automates enrichment from an arbitrary input prompt.

  \item \textbf{No integrated system.} No prior work combines enrichment,
        scanning, and repair in a single runnable tool. The components exist
        in separate papers; they have never been assembled.

  \item \textbf{No reproducible multi-model dataset.} No prior study logs all
        runs---prompts, enriched prompts, raw outputs, scan findings, repair
        iterations---to a queryable database for subsequent analysis.
\end{enumerate}

LeBlanc addresses all five. The system combines automated enrichment (closing
gap~3), multi-model parallel evaluation (closing gap~1), iterative automated
repair with re-scanning (closing gap~2), an integrated runnable pipeline
(closing gap~4), and a logged SQLite dataset (closing gap~5).

\printbibliography[heading=none]

\end{document}
